\section{The Reinforcement Learning Paradigm}

The main goal of \gls{rl} is to conceive of an \emph{agent} that can learn to maximise some notion of utility by interacting with its \emph{environment}. 
The interaction is modelled as a sequential decision problem. 
At each step, the agent can apply an \emph{action} in the environment.
As a consequence of this action and external factors, the environment changes to a new \emph{state} and provides the agent with a reinforcement signal.
This signal, called \emph{reward}, gives feedback on the quality of the action. 
The agent's goal is to maximise the cumulative sum of rewards-- or \emph{return}.
The sequential process followed by the environment, involving transition dynamics and rewards, is modelled by a \glsfirst{mdp}. 
To maximise its rewards, the agent must autonomously balance between exploration and exploitation.  
It must explore new areas of the \gls{mdp} to avoid missing some opportunities but should also consider exploiting known profitable parts of the \gls{mdp} to maximise its utility. 
The framework of \gls{rl} is inspired by discoveries from psychology on the human learning process. 
It is general enough to include many problems and in particular the other subfields of \gls{ml} that are \emph{supervised} and \emph{unsupervised learning}. 
It is also specific enough to yield meaningful theoretical results that guide the design of algorithms.

Empirically, the soundness of the framework is supported by the breakthrough that it has allowed. 
A limited sample of these achievements is:
robotic locomotion \cite{haarnoja2018learning};
autonomous driving \cite{kiran2021deep};
playing complex video games at master level \cite{vinyals2019grandmaster}; designing computer chips that perform more and consume less energy \cite{mirhoseini2021graph};
controlling the shape of a plasma in a tokamak \cite{degrave2022magnetic}.


However, most of these applications require that the original \gls{rl} framework be slightly extended to account for partial observability of the state, for example. 
Another limitation of the original framework could be the presence of delay in the dynamical system. 
In the following section, we explain why taking delay into account is the key to achieving great empirical results.
For so fast-paced \gls{ml} and \gls{rl} research can be, we propose to the reader to take the time to delve into the problem of delays.
